% ============================================================
%  CHAPITRE 1 — INTRODUCTION
%  Présentation TalentMatch — 1ère Restitution
% ============================================================
\documentclass[aspectratio=169]{beamer}

% ── Thème ────────────────────────────────────────────────────
\usetheme{Madrid}
\usecolortheme{default}
\setbeamercolor{palette primary}{bg=blue!70!black, fg=white}
\setbeamercolor{palette secondary}{bg=blue!60!black, fg=white}
\setbeamercolor{palette tertiary}{bg=blue!50!black, fg=white}
\setbeamercolor{palette quaternary}{bg=blue!40!black, fg=white}
\setbeamercolor{structure}{fg=blue!70!black}
\setbeamercolor{block title}{bg=blue!70!black, fg=white}
\setbeamercolor{block body}{bg=blue!5!white}
\setbeamercolor{alertblock title}{bg=red!70!black, fg=white}
\setbeamercolor{exampleblock title}{bg=green!50!black, fg=white}

% ── Packages ─────────────────────────────────────────────────
\usepackage[utf8]{inputenc}
\usepackage[T1]{fontenc}
\usepackage[french]{babel}
\usepackage{graphicx}
\usepackage{booktabs}
\usepackage{tikz}
\usetikzlibrary{shapes.geometric, arrows.meta, positioning}
\usepackage{fontawesome5}
\usepackage{hyperref}

% ── Métadonnées ───────────────────────────────────────────────
\title[TalentMatch — 1\`ere Restitution]{%
  \textbf{TalentMatch}\\
  \large Plateforme intelligente de matching CV / Offres d'emploi}
\subtitle{Rapport de stage — 1\`ere Restitution}
\author[Youssef Gara]{%
  \textbf{Youssef Gara}\\[4pt]
  {\small Encadrantes : Mme Loujayne Bouzarti \& Mme Sirine Nasri}}
\institute[3S]{%
  \textbf{3S}\\
  Stage du \textbf{Février 2026} au \textbf{Août 2026}}
\date{\today}

% ── Pied de page personnalisé ─────────────────────────────────
\setbeamertemplate{footline}{%
  \leavevmode%
  \hbox{%
    \begin{beamercolorbox}[wd=.33\paperwidth,ht=2.5ex,dp=1ex,center]{palette secondary}
      \usebeamerfont{author in head/foot}\insertshortauthor
    \end{beamercolorbox}%
    \begin{beamercolorbox}[wd=.34\paperwidth,ht=2.5ex,dp=1ex,center]{palette tertiary}
      \usebeamerfont{title in head/foot}\insertshorttitle
    \end{beamercolorbox}%
    \begin{beamercolorbox}[wd=.33\paperwidth,ht=2.5ex,dp=1ex,right]{palette primary}
      \usebeamerfont{date in head/foot}
      \insertframenumber{} / \inserttotalframenumber\hspace*{2ex}
    \end{beamercolorbox}}%
  \vskip0pt%
}
\setbeamertemplate{navigation symbols}{}

% ============================================================
\begin{document}
% ============================================================

% ── SLIDE TITRE ──────────────────────────────────────────────
\begin{frame}
  \titlepage
\end{frame}

% ── SLIDE PLAN GÉNÉRAL ───────────────────────────────────────
\begin{frame}{Plan de la présentation}
  \tableofcontents[hideallsubsections]
\end{frame}

% ============================================================
\section{Introduction}
% ============================================================

% ── SLIDE 1 : CONTEXTE DU MARCHÉ ─────────────────────────────
\begin{frame}{Contexte — Le défi du recrutement moderne}

  \begin{columns}[T]
    \begin{column}{0.48\textwidth}
      \begin{alertblock}{\faExclamationTriangle\ \ Problèmes actuels}
        \begin{itemize}
          \item Un recruteur consulte un CV en moyenne \textbf{23 secondes}
          \item Les outils existants filtrent par \textbf{mots-clés uniquement}
          \item Pas de compréhension du \textbf{sens réel} des profils
          \item Nombreux \textbf{faux positifs} et \textbf{bons profils ignorés}
        \end{itemize}
      \end{alertblock}
    \end{column}
    \begin{column}{0.48\textwidth}
      \begin{exampleblock}{\faLightbulb\ \ Notre réponse}
        \begin{itemize}
          \item Analyse sémantique par \textbf{modèles IA} (NLP)
          \item Compréhension du \textbf{domaine professionnel} (ISCO-08)
          \item Score de matching \textbf{explicable} et \textbf{calibré}
          \item Classement automatique des candidats
        \end{itemize}
      \end{exampleblock}
    \end{column}
  \end{columns}

  \vspace{0.5cm}
  \begin{block}{}
    \centering
    \large
    \textit{Comment faire correspondre \textbf{intelligemment} un candidat à une offre,\\
    au-delà de la simple recherche par mots-clés ?}
  \end{block}

\end{frame}

% ── SLIDE 2 : OBJECTIF DU STAGE ──────────────────────────────
\begin{frame}{Objectif du stage}

  \begin{block}{\faBullseye\ \ Mission principale}
    Concevoir et développer \textbf{TalentMatch} : une plateforme web complète
    de matching automatique CV / Offres d'emploi, intégrant des modèles
    d'intelligence artificielle et des algorithmes de NLP avancés.
  \end{block}

  \vspace{0.4cm}

  \begin{columns}[T]
    \begin{column}{0.23\textwidth}
      \centering
      \textcolor{blue!70!black}{\faFilePdf}\\[4pt]
      \textbf{Extraction}\\
      \small Analyser et structurer les CVs automatiquement
    \end{column}
    \begin{column}{0.23\textwidth}
      \centering
      \textcolor{blue!70!black}{\faBrain}\\[4pt]
      \textbf{Compréhension}\\
      \small Comprendre la sémantique des profils par IA
    \end{column}
    \begin{column}{0.23\textwidth}
      \centering
      \textcolor{blue!70!black}{\faChartBar}\\[4pt]
      \textbf{Scoring}\\
      \small Calculer un score de pertinence calibré
    \end{column}
    \begin{column}{0.23\textwidth}
      \centering
      \textcolor{blue!70!black}{\faListOl}\\[4pt]
      \textbf{Classement}\\
      \small Classer et expliquer les résultats
    \end{column}
  \end{columns}

  \vspace{0.5cm}
  \begin{columns}[T]
    \begin{column}{0.48\textwidth}
      \textbf{Cadre :}
      \begin{itemize}
        \item Entreprise : \textbf{3S}
        \item Encadrantes : Mme \textbf{Loujayne Bouzarti} \& Mme \textbf{Sirine Nasri}
      \end{itemize}
    \end{column}
    \begin{column}{0.48\textwidth}
      \textbf{Période :}
      \begin{itemize}
        \item Stage de \textbf{Février 2026} à \textbf{Août 2026}
        \item Durée : \textbf{6 mois}
      \end{itemize}
    \end{column}
  \end{columns}

\end{frame}

% ── SLIDE 3 : VUE D'ENSEMBLE DE LA SOLUTION ──────────────────
\begin{frame}{TalentMatch — Vue d'ensemble}

  \begin{center}
  \begin{tikzpicture}[
    node distance=0.7cm and 1.2cm,
    box/.style={rectangle, rounded corners=4pt, draw=blue!60!black,
                fill=blue!10, text width=2.4cm, align=center,
                minimum height=1cm, font=\small},
    arrow/.style={-{Stealth[length=4pt]}, thick, blue!60!black},
  ]
    \node[box] (cv)      {\faFilePdf\\[2pt]\textbf{CV (PDF)}};
    \node[box, right=of cv] (ext)   {\faCogs\\[2pt]\textbf{Extraction}\\NLP};
    \node[box, right=of ext] (match) {\faBrain\\[2pt]\textbf{Scoring}\\IA (BGE-M3)};
    \node[box, right=of match] (rank) {\faListOl\\[2pt]\textbf{Classement}\\Candidats};
    \node[box, below=of match] (offer) {\faBriefcase\\[2pt]\textbf{Offre}\\d'emploi};

    \draw[arrow] (cv)    -- (ext);
    \draw[arrow] (ext)   -- (match);
    \draw[arrow] (match) -- (rank);
    \draw[arrow] (offer) -- (match);
  \end{tikzpicture}
  \end{center}

  \vspace{0.3cm}
  \begin{columns}[T]
    \begin{column}{0.32\textwidth}
      \begin{block}{\small Backend}
        \centering\small FastAPI · Python\\spaCy · PostgreSQL
      \end{block}
    \end{column}
    \begin{column}{0.32\textwidth}
      \begin{block}{\small Modèles IA}
        \centering\small BGE-M3 · Cross-Encoder\\ISCO-08 (53 domaines)
      \end{block}
    \end{column}
    \begin{column}{0.32\textwidth}
      \begin{block}{\small Frontend}
        \centering\small React 18 · Vite\\Tailwind CSS
      \end{block}
    \end{column}
  \end{columns}

\end{frame}

% ── SLIDE 4 : PLAN DE LA PRÉSENTATION ────────────────────────
\begin{frame}{Plan de cette présentation}
  \begin{enumerate}
    \setlength\itemsep{0.5em}
    \item[\textbf{1.}] \textbf{Introduction} \hfill \textcolor{gray}{\small ← vous êtes ici}
    \item[\textbf{2.}] Cadre du projet \textit{(problématique, existant, solution)}
    \item[\textbf{3.}] Analyse des besoins \textit{(fonctionnels \& non fonctionnels)}
    \item[\textbf{4.}] Architecture \textit{(backend, frontend, IA)}
    \item[\textbf{5.}] Méthodologie \textit{(Agile, outils, organisation)}
    \item[\textbf{6.}] Avancement actuel \textit{(démos, résultats)}
    \item[\textbf{7.}] Conclusion \& Perspectives
  \end{enumerate}
\end{frame}

% ============================================================
\end{document}
% ============================================================
